\par\medskip
\noindent
2.5. El teorema de ortogonalidad posee una formulación equivalente
mediante una polarización $\xi$ de $A_K$ y la forma bilineal alternante
(1.0.8), compatible con las operaciones de $\pi$,
\begin{equation}\tag{2.5.1}
 \varphi_\xi:T_\ell(A(\overline K))\times T_\ell(A(\overline K))
 \longrightarrow T_\ell(\overline K^*),
\end{equation}
obtenida de la forma $\varphi$ de (1.0.3) transportando la estructura
en el segundo argumento a lo largo de la isogenia de
$\mathbb Z_\ell$-módulos
\[
 T_\ell(A(\overline K))\longrightarrow T_\ell(A'(\overline K))
\]
inducida por la isogenia $\psi_\xi:A_K\longrightarrow A'_K$ definida
por $\xi$. La igualdad de (2.4.6) adopta entonces la forma sencilla
\begin{equation}\tag{2.5.2}
 W=V\cap V^\perp,
\end{equation}
donde hemos puesto
\begin{equation}\tag{2.5.3}
 \left\{
 \begin{array}{l}
 U=T_\ell(A(\overline K)),\qquad
 V=U^I=\text{la «parte fija» de }T_\ell(A(\overline K)),\\[2mm]
 \hspace{4em}W=\text{la parte tórica de }T_\ell(A(\overline K)),
 \end{array}
 \right.
\end{equation}
y $\perp$ denota el ortogonal respecto de la forma de polarización
(2.5.1). Estos módulos dan el siguiente diagrama de inclusiones:
\begin{equation}\tag{2.5.4}
\begin{tikzcd}[column sep=large,row sep=large]
 && V\arrow[dl,hook',"2\lambda"'] && \\
 U & W^\perp=V+V^\perp\arrow[l,hook',"\mu"']
 && W=V\cap V^\perp
    \arrow[ul,hook',"2\alpha"']
    \arrow[dl,hook,"2\lambda"]
 & 0\arrow[l,hook',"\mu"'] \\
 && V^\perp\arrow[ul,hook,"2\alpha"] &&
\end{tikzcd}
\qquad ,
\end{equation}
donde los enteros $2\alpha$, $2\lambda$ y $\mu$ son los rangos de los
cocientes de módulos consecutivos en el diagrama, con la notación de
(2.1.9). (Que $W$ tenga rango $\mu$ es la definición de $\mu$; que
$V/W$ tenga rango $2\alpha$ se deduce de la segunda relación de
(2.1.11), junto con la sucesión exacta (2.1.6). Por tanto, por
ortogonalidad, $\operatorname{rank}(U/W^\perp)=\mu$ y
$\operatorname{rank}(W^\perp/V^\perp)=2\alpha$. Finalmente, las
fórmulas
\begin{equation}\tag{2.5.5}
 \operatorname{rank}\bigl(\sisi{V}{V^\perp}/(V\cap V^\perp)\bigr)
 =\operatorname{rank}(\sisi{V}{V^\perp}/W)=2\lambda
\end{equation}
y la fórmula dual se deducen de las precedentes, en vista de (2.1.10)
y de $\operatorname{rank}U=2n$.)

\medskip
\noindent
\underline{Observaciones} 2.6.\quad a) Hemos obtenido aquí el teorema
de ortogonalidad 2.4 como consecuencia fácil de la teoría de
biextensiones desarrollada en las secciones precedentes. La
\underline{inclusión} (2.4.6) puede demostrarse sin utilizar la teoría
de biextensiones, mediante un argumento aritmético al estilo de la
demostración del «teorema de monodromía» dada en el Exposé I. Se reduce
a una situación en la que $S$ posee cuerpo residual finito, donde el
resultado de ortogonalidad se deduce de una consideración sencilla de
los \underline{pesos} de la operación de Frobenius (utilizando el
resultado de A.~Weil sobre los valores propios de este Frobenius). No
es necesario invocar aquí la resolución de singularidades: puede
enunciarse el resultado deseado para esquemas en grupos $A$ sobre
esquemas $S$ bajo condiciones suficientemente generales. Por desgracia,
no parece posible obtener por este método la igualdad en la inclusión
precedente. Al parecer, esta igualdad se utilizará de manera esencial
en la demostración del «teorema de reducción semiestable» de la sección
siguiente.

\smallskip
\noindent
b) El teorema de ortogonalidad 2.4 permite expresar las propiedades más
importantes del modelo de Néron de una variedad abeliana $A_K$ en
términos de la representación del grupo de inercia $I$ sobre el módulo
de Tate $T_\ell(A(\overline K))$, donde $\ell$ es un primo fijo distinto
de la característica residual. (El caso $\ell=p$ se considerará por
separado más adelante, en la sección 5.) Como primer enunciado de este
tipo, aunque no utiliza el teorema de ortogonalidad, recordemos el
criterio de buena reducción 2.2.9. Daremos a continuación muchos otros
ejemplos, comenzando en la sección siguiente (3.5).
